\documentclass[11pt, a4paper]{article}

\usepackage[utf8]{inputenc}
\usepackage[T1]{fontenc}
\usepackage[margin=2.5cm]{geometry}
\usepackage{booktabs}
\usepackage{xcolor}
\usepackage{enumitem}
\usepackage{titlesec}
\usepackage{hyperref}
\usepackage{microtype}
\usepackage{parskip}

% ── Colours ──────────────────────────────────────────────────────────────────
\definecolor{claimblue}{RGB}{30, 100, 180}
\definecolor{justgreen}{RGB}{20, 130, 60}
\definecolor{warnred}{RGB}{180, 40, 40}
\definecolor{lightgray}{RGB}{245, 245, 245}

% ── Section style ─────────────────────────────────────────────────────────────
\titleformat{\section}{\large\bfseries\color{claimblue}}{}{0em}{}[\titlerule]
\titleformat{\subsection}{\normalsize\bfseries}{}{0em}{}

% ── Custom commands ───────────────────────────────────────────────────────────
\newcommand{\Claim}[1]{%
  \vspace{6pt}%
  \noindent\fcolorbox{claimblue}{lightgray}{\parbox{\dimexpr\linewidth-2\fboxsep-2\fboxrule}{%
    \textbf{\textcolor{claimblue}{Claim:}} #1}}%
  \par\vspace{4pt}%
}

\newcommand{\Why}[1]{%
  \noindent\fcolorbox{justgreen}{white}{\parbox{\dimexpr\linewidth-2\fboxsep-2\fboxrule}{%
    \textbf{\textcolor{justgreen}{Why / How to explain it:}} #1}}%
  \par\vspace{8pt}%
}

\newcommand{\Numbers}[1]{%
  \noindent\textit{\textcolor{gray}{Numbers: #1}}\par\vspace{2pt}%
}

% ─────────────────────────────────────────────────────────────────────────────
\begin{document}

\begin{center}
  {\LARGE\bfseries Claims \& Oral Justifications}\\[4pt]
  {\large How Unified is your Encoder? --- Results \& Conclusion}\\[2pt]
  {\small CS-503 Visual Intelligence --- Final Presentation --- May 2026}
\end{center}

\vspace{8pt}
\noindent\textit{This document lists every factual claim made in the Results and Conclusion slides,
with a plain-English justification that is easy to deliver orally.
Numbers come directly from the experimental outputs.}

\vspace{8pt}\hrule\vspace{8pt}

% =============================================================================
\section{RQ1 --- Cross-Modal Similarity (Separate Passes)}
% =============================================================================

\subsection{Result slide 1}

\Claim{Cross-modal similarity rises monotonically across encoder layers for all three metric types (CKA, PWCCA, k-NN) and all three modality pairs.}

\Numbers{%
  Hypersim, CKA: rgb--depth $0.06 \to 0.40$; depth--normals $0.13 \to 0.33$; rgb--normals $0.03 \to 0.20$.
  Pattern is consistent across DIODE.%
}

\Why{%
  We fed each modality separately into the encoder, layer by layer.
  At the very first layer the representations are nearly orthogonal --- the encoder
  has not had time to mix them yet, so CKA is close to 0.
  As we go deeper, the encoder processes the tokens through 12 self-attention blocks
  and gradually ``pulls'' them toward a common geometric structure.
  By layer 11, all three pairs are far above the null (see next claim).
  The fact that \emph{three independent metrics} (which capture geometry, direction, and topology)
  all agree rules out any single-metric artefact.%
}

\subsection{Result slide 2}

\Claim{The signal is highly statistically significant: p-value $= 0.004$ (minimum across layers and pairs).}

\Numbers{%
  Null mean at layer 11, CKA: $\approx 0.005$ (all pairs).
  Matched CKA at layer 11: $0.20$--$0.40$ depending on pair.
  That is $40$--$80\times$ above the null mean.%
}

\Why{%
  We built a null distribution by comparing \emph{mismatched} scenes ---
  taking the RGB of scene A and the depth of a randomly chosen, different scene B,
  making sure the two scenes belong to different room types (e.g.\ bathroom vs.\ office).
  This null captures any similarity that is purely architectural (shared positional encodings,
  layer norms) rather than learned.
  The matched signal sits $40$--$80\times$ above this null at deep layers.
  We repeated the shuffle 1\,000 times to get a stable estimate.
  The reported p-value of 0.004 is the \emph{most conservative} value across all pairs and layers,
  so the claim holds at every point on the plot.%
}

\subsection{Result slide 3}

\Claim{DIODE shows uniformly higher similarity than Hypersim, likely due to DIODE's lower visual diversity.}

\Why{%
  Hypersim contains 100 varied synthetic scenes (offices, kitchens, hallways\ldots).
  DIODE's indoor/outdoor split covers fewer scene types.
  When there is less variety in the dataset, the scene representations from any two modalities
  naturally land closer together in latent space --- even before the encoder does any work.
  This raises the floor of similarity, shifting the entire curve up.
  Importantly, \emph{both} datasets show the same monotonic rise and the same significance,
  so the conclusion about convergence is not an artefact of one dataset.%
}

\vspace{8pt}\hrule\vspace{8pt}

% =============================================================================
\section{RQ1b --- Joint Pass vs Separate Pass}
% =============================================================================

\Claim{Joint-pass similarity is uniformly higher than separate-pass similarity.}

\Why{%
  In the \emph{separate} pass, each modality goes through the encoder on its own.
  In the \emph{joint} pass, both modalities are concatenated into a single token sequence
  and processed together.
  When the tokens can attend to each other through self-attention, the encoder
  explicitly mixes them --- there is cross-modal attention happening inside every block.
  This is not the case when they are processed separately.
  The result is that the joint representations are consistently more similar: the encoder
  is doing more ``merging work'' when it sees both modalities at once.%
}

\Claim{The gain concentrates at deep layers and is largest for the hardest pair: RGB $\leftrightarrow$ Normals.}

\Why{%
  At early layers, joint and separate curves are nearly identical --- the tokens have not had
  time to interact yet.
  As depth increases, cross-modal attention accumulates and the gap opens up.
  RGB and surface normals are the most \emph{semantically distant} pair:
  RGB encodes color and texture, normals encode local surface orientation.
  These two streams have the most to ``learn from each other,'' so when they can attend
  to each other (joint pass), the gain is the largest.%
}

\vspace{8pt}\hrule\vspace{8pt}

% =============================================================================
\section{RQ2 --- Transmodal Triangulation (CKA Fragmentation)}
% =============================================================================

\Claim{At layer 11, $\geq 86.5\%$ of cross-modal geometric structure is linearly unified (average across pairs).}

\Numbers{%
  Depth--RGB: $91.4\%$ unified (fragmentation $= 8.6\%$).
  Normals--RGB: $87.5\%$ unified (fragmentation $= 12.5\%$).
  Depth--Normals: $80.6\%$ unified (fragmentation $= 19.4\%$).
  Average: $86.5\%$.%
}

\Why{%
  The CKA fragmentation score measures how much of the shared structure between modality A and
  modality B \emph{cannot} be linearly explained through a third modality C acting as a bridge.
  Concretely: we project A and B through C and ask ``how much of what A and B share is also
  accessible from C?''
  A score of 0 means perfect triangulation (C perfectly mediates the A--B relationship).
  A score of 1 means total fragmentation (C adds nothing).
  Getting an average of $86.5\%$ unified means that almost nine tenths of the cross-modal
  geometric structure is linearly shared --- even when measured through a third modality.%
}

\Claim{RGB-involving pairs unify strongly; Depth--Normals converges much more slowly, suggesting RGB acts as a geometric hub.}

\Numbers{%
  Depth--RGB fragmentation: $35.6\% \to 8.6\%$ (drop of 27 pp across layers).
  Normals--RGB fragmentation: $30.3\% \to 12.5\%$ (drop of 18 pp).
  Depth--Normals fragmentation: $22.6\% \to 19.4\%$ (drop of only 3 pp).%
}

\Why{%
  The two pairs that involve RGB converge strongly and consistently across all 12 layers.
  The pair that does \emph{not} involve RGB (Depth--Normals) barely moves: it starts at
  $22.6\%$ fragmentation and ends at $19.4\%$ --- essentially flat.
  One natural interpretation: the encoder uses the RGB subspace as a common coordinate system
  that both Depth and Normals get mapped into.
  Depth and Normals then share structure \emph{indirectly}, via RGB.
  Without RGB as intermediary, their representations do not converge further.
  This is consistent with RGB being the dominant modality in 4M training
  (mono-modal RGB inputs are far more frequent than mono-modal depth or normals inputs).%
}

\Claim{Residual fragmentation $> 0$ at all layers and all pairs: 4M is measurably not perfectly unified ($p < 0.001$, BH-FDR corrected).}

\Numbers{%
  All six triplet tests at layer 11 have $p_{\text{adj}} = 0.001$ (minimum detectable with 1\,000
  bootstrap draws).
  The minimum residual fragmentation at layer 11 is $8.6\%$ (Depth--RGB).%
}

\Why{%
  Even in the best case, $8.6\%$ of the cross-modal structure is \emph{not} linearly unified.
  This is not noise: we tested it against a bootstrap null (shuffled triplets) and the result
  is significant at $p < 0.001$ after Benjamini--Hochberg correction for multiple comparisons.
  In plain terms: the encoder is very good at converging, but it is not collapsing all modalities
  onto a single shared representation.
  Some modality-specific structure remains, encoded in subspaces that the other modalities
  cannot access.%
}

\vspace{8pt}\hrule\vspace{8pt}

% =============================================================================
\section{Conclusion}
% =============================================================================

\Claim{Yes, the encoder converges: it is not a hidden ensemble.}

\Why{%
  Zamir's conjecture was that 4M might secretly act as an ensemble of independent sub-networks,
  one per modality.
  Our results reject this: cross-modal similarity is $40$--$80\times$ above the null at deep layers,
  consistently across three metrics, two datasets, and three modality pairs.
  An ensemble of independent sub-networks would produce similarity close to the null --- we see the
  exact opposite.%
}

\Claim{Joint processing strengthens unification, especially at deep layers.}

\Why{%
  When modalities are fed together, the encoder does more cross-modal mixing via self-attention.
  The effect is strongest at deep layers (where many attention blocks have accumulated) and
  for the hardest pair (RGB--Normals), which has the most to gain from explicit interaction.%
}

\Claim{One trunk, but imperfect: 4M converges, but does not fully unify ($\geq 8.6\%$ residual fragmentation, $p < 0.001$).}

\Why{%
  We answered the question ``to what degree'' --- not just yes/no.
  $86.5\%$ of cross-modal structure is linearly unified at the last encoder layer.
  But the remaining $13.5\%$ is statistically proven to be real fragmentation, not noise.
  The encoder is one trunk --- but not a perfectly seamless one.
  This is consistent with the Platonic Representation Hypothesis (representations converge
  across modalities and models) while adding a quantitative bound on how far 4M has gone.%
}

\vspace{8pt}\hrule\vspace{8pt}

% =============================================================================
\section{Exploratory Analysis --- t-SNE Slide (Slide 5)}
% =============================================================================

\subsection{What we observe}

The t-SNE plot shows DIODE encoder activations at each layer, coloured in three ways:
(1) by modality (RGB vs.\ Depth), (2) by scene category, (3) by scene id.

The most striking observation is that \textbf{two distinct, well-separated clusters appear when coloured by modality} --- one for RGB, one for Depth --- and they do \emph{not} merge as layer depth increases.

\subsection{Does this contradict our conclusion?}

\Claim{Two persistent modality clusters in t-SNE are \emph{consistent} with high CKA convergence --- they measure different things.}

\Why{%
  \textbf{t-SNE shows \emph{where} representations live in space.}
  CKA measures \emph{how they are organised relative to each other} (their pairwise distance structure).
  These two questions are independent.

  \medskip
  \textbf{Analogy:} imagine two maps of the same city --- one labelled in French (RGB) and one in English (Depth).
  Overlaid on a single plot, you would see two clusters of words.
  But the \emph{street layout} (relational structure) is identical in both maps.
  CKA measures the street layout. t-SNE measures where the word labels land.

  \medskip
  This phenomenon is known as the \textbf{modality gap} (Liang et al., NeurIPS 2022):
  in multimodal models, different modalities systematically occupy distinct regions of the
  latent space even when they encode similar information.
  CKA is invariant to this global offset --- it only cares about whether the
  pairwise distances \emph{between scenes} are the same across modalities.%
}

\Claim{The persistent cluster separation actually \emph{supports} the ``one trunk, but imperfect'' conclusion.}

\Why{%
  If the encoder were perfectly unified, all modality-specific structure would vanish and the
  two clouds would collapse into one.
  The fact that they do \emph{not} merge is the geometric signature of the residual fragmentation
  we quantify in RQ2: $\geq 8.6\%$ of the cross-modal structure remains fragmented at layer 11
  ($p < 0.001$).
  The t-SNE is therefore not a contradiction --- it is a \emph{visual confirmation} of the
  same imperfection that our fragmentation score measures numerically.%
}

\Claim{The \emph{within-cluster} structure (scene id colouring) evolves with layer depth and is consistent across modalities --- this is what CKA captures.}

\Why{%
  If you look at the third panel (coloured by scene id), scenes that are similar tend to land
  in similar relative positions \emph{within} their respective modality cluster,
  and this organisation becomes cleaner at deeper layers.
  The two clouds stay apart, but the internal geometry of each cloud converges toward the same pattern.
  CKA measures exactly this: it asks whether the matrix of pairwise scene distances looks the
  same for RGB activations as for Depth activations.
  High CKA means yes --- even if the two clouds are nowhere near each other in absolute position.%
}

\subsection{How to say it at the oral}

\begin{quote}
\itshape
``You might wonder: if the encoder converges, why do we still see two distinct clusters here?
The answer is that t-SNE and CKA measure different things.
t-SNE shows \emph{where} representations live in space --- and yes, each modality occupies its own region.
CKA measures whether the \emph{relational structure between scenes} is the same across modalities.
Two clouds can be in different places but be organised identically internally.
The persistent separation is actually consistent with our RQ2 result: there is always some
residual modality-specific structure, which is exactly what we call residual fragmentation.''
\end{quote}

\vspace{8pt}\hrule\vspace{8pt}

% =============================================================================
\section*{Quick reference: key numbers at a glance}
% =============================================================================

\begin{center}
\begin{tabular}{lcc}
\toprule
\textbf{Claim} & \textbf{Value} & \textbf{Evidence} \\
\midrule
Min p-value (RQ1, all pairs/layers)            & $0.004$           & 1\,000-draw null, Hypersim+DIODE \\
CKA at layer 11, rgb--depth (Hypersim)         & $0.40$            & $66\times$ above null mean \\
Null mean at layer 11, CKA                     & $0.005$           & 1\,000 mismatched draws \\
Average unification at layer 11 (RQ2)          & $86.5\%$          & CKA fragmentation score \\
Best pair (Depth--RGB) unification             & $91.4\%$          & fragmentation $= 8.6\%$ \\
Worst pair (Depth--Normals) unification        & $80.6\%$          & fragmentation $= 19.4\%$ \\
Residual fragmentation p-value (RQ2, layer 11) & $< 0.001$         & BH-FDR, 1\,000 bootstrap draws \\
Depth--Normals fragmentation drop (L0$\to$L11) & $22.6\% \to 19.4\%$ & only $-3$ pp (vs.\ $-27$ pp for Depth--RGB) \\
\bottomrule
\end{tabular}
\end{center}

\end{document}
